\section{Resultados}

El HP con nombre \textit{Tall stature} y código HP:0000098 tiene 182 genes que al mapearlos en la red de STRING \cite{Szklarczyk2023} con un umbral de corte de 900, quedan 166 genes y 300 conexiones. La red resultante tiene una transitivity de 0,5016.

El HP con nombre \textit{Tall stature} y código HP:0000098 tiene 182 genes que al mapearlos en la red de STRING \cite{Szklarczyk2023} con un umbral de corte de 900, quedan 166 genes y 300 conexiones. La red resultante tiene una transitivity de 0,5016.

El HP con nombre \textit{Tall stature} y código HP:0000098 tiene 182 genes que al mapearlos en la red de STRING \cite{Szklarczyk2023} con un umbral de corte de 900, quedan 166 genes y 300 conexiones. La red resultante tiene una transitivity de 0,5016.

El HP con nombre \textit{Aortic dissection} y código HP:0002647 tiene 41 genes que al mapearlos en la red de STRING \cite{Szklarczyk2023} con un umbral de corte de 700, quedan 39 genes y 158 conexiones (79 interacciones, ya que STRING incluye cada par dos veces). La red resultante tiene una transitivity de 0,5758.

\subsection{Propiedades de la red}
La red de interacción de STRING (versión 12.5, umbral 700) de los genes asociados a HP:0002647 se construyó con todos los genes encontrados en STRING, de modo que los genes sin ninguna interacción quedaron como nodos aislados. La Tabla~\ref{tab:red} resume sus propiedades (archivo \texttt{results/minimo/tablas/topologia\_resumen.tsv}). Diez de los 41 genes quedaron aislados: \textit{SLC2A10}, \textit{IPO8}, \textit{FOXE3}, \textit{PRKG1}, \textit{MAT2A}, \textit{THSD4}, \textit{ANGPTL6}, \textit{THSD1}, \textit{ROBO4} y \textit{GATA5}. Los genes con más interacciones fueron \textit{TGFBR1} y \textit{TGFBR2} (11 interacciones cada uno) y \textit{FBN1} (10) (archivo \texttt{topologia\_genes.tsv}).

\begin{table}[h!]
	\renewcommand{\arraystretch}{1.25}
	\caption{Propiedades de la red de interacción de los genes asociados a HP:0002647 (STRING 12.5, umbral 700). Fuente: \texttt{results/minimo/tablas/topologia\_resumen.tsv}.}
	\label{tab:red}
	\begin{tabular}{lrp{6.3cm}}
		\hline
		Propiedad & Valor & Interpretación \tabularnewline
		\hline
		Nodos & 41 & Todos los genes de HPO se encontraron en STRING \tabularnewline
		Nodos con interacción & 31 & Genes con al menos una interacción con otro gen de la lista \tabularnewline
		Nodos aislados & 10 & Genes sin ninguna interacción con otro gen de la lista \tabularnewline
		Aristas & 75 & Interacciones entre pares de genes \tabularnewline
		Grado medio & 3,659 & Cada gen interacciona, de media, con unos 3,7 genes de la lista \tabularnewline
		Densidad & 0,0915 & Proporción de las 820 parejas posibles que están conectadas \tabularnewline
		Componentes conexas & 13 & La red no es conexa: tiene 13 bloques, entre ellos los 10 genes aislados \tabularnewline
		Nodos de la componente mayor & 26 & 26 de los 31 genes conectados forman un único bloque \tabularnewline
		\hline
	\end{tabular}
\end{table}